\newcommand{\id}{50}
\newcommand{\nom}{Plan des moindres carrés}
\newcommand{\sequence}{04}
\newcommand{\numero}{16}
\newcommand{\type}{TD}
\newcommand{\descrip}{Déterminer le plan des moindres carrés à partir d'un nuage de points.}

\documentclass[10pt,a4paper]{article}
  \usepackage[french]{babel}
  \usepackage[utf8]{inputenc}
  \usepackage[T1]{fontenc}
  \usepackage{xcolor}
  \usepackage[]{graphicx}
  \usepackage{makeidx}
  \usepackage{textcomp}
  \usepackage{amsmath}
  \usepackage{amssymb}
  \usepackage{stmaryrd}
  \usepackage{fancyhdr}
  \usepackage{lettrine}
  \usepackage{calc}
  \usepackage{boxedminipage}
  \usepackage[french,onelanguage, boxruled,linesnumbered]{algorithm2e}
  \usepackage[colorlinks=false,pdftex]{hyperref}
  \usepackage{minted}
  \usepackage{url}
  %\usepackage[locale=FR]{siunitx}
  \usepackage{multicol}
  \makeindex

  %\graphicspath{{../Images/}}

  \renewcommand\listingscaption{Programme}

  %\renewcommand{\thechapter}{\Alph{chapter}}
  \renewcommand{\thesection}{\Roman{section}}
  %\newcommand{\inter}{\vspace{0.5cm}%
  %\noindent }
  %\newcommand{\unite}{\ \textrm}
  \newcommand{\ud}{\mathrm{d}}
  \newcommand{\vect}{\overrightarrow}
  %\newcommand{\ch}{\mathrm{ch}} % cosinus hyperbolique
  %\newcommand{\sh}{\mathrm{sh}} % sinus hyperbolique

  \textwidth 160mm
  \textheight 250mm
  \hoffset=-1.70cm
  \voffset=-1.5cm
  \parindent=0cm

  \pagestyle{fancy}
  \fancyhead[L]{\bfseries {\large PTSI -- Dorian}}
  \fancyhead[C]{\bfseries{{\type} \no \num}}
  \fancyhead[R]{\bfseries{\large Informatique}}
  \fancyfoot[C]{\thepage}
  \fancyfoot[L]{\footnotesize R. Costadoat, J. Genzmer, W. Robert}
  \fancyfoot[R]{\small \today}
  
  \definecolor{bg}{rgb}{0.5,0.5,0.5}
  \definecolor{danger}{RGB}{217,83,79}
  
  \fancypagestyle{correction}{%
  \fancyhf{}
  \lhead{\colorbox{danger}{\begin{minipage}{0.65\paperwidth} \textcolor{white}{\textbf{Correction}} \end{minipage}} }
  \rhead{\includegraphics[width=2cm]{../../img/logo}}
  \lfoot{Juliette Genzmer, Willie Robert, Renaud Costadoat}
  \rfoot{\colorbox{danger}{\begin{minipage}{0.6\paperwidth} \begin{flushright}\textcolor{white}{\textbf{Correction}}\end{flushright} \end{minipage}} }}

  
  % macro Juliette
  
\usepackage{comment}   
\usepackage{amsthm}  
\theoremstyle{definition}
\newtheorem{exercice}{Exercice}
\newtheorem*{rappel}{Rappel}
\newtheorem*{remark}{Remarque}
\newtheorem*{defn}{Définition}
\newtheorem*{ppe}{Propriété}
\newtheorem{solution}{Solution}

 
\begin{document}

\section{Position du problème}

On dispose d'un nuage de $n$ points $(x_i, y_i, z_i)$ pour $i = 1, \dots, n$.
On souhaite trouver le plan d'équation
\[
z = a x + b y + c
\]
qui s'ajuste « au mieux » à ces points, au sens des moindres carrés.

\section{Fonction d'erreur}

Pour un point $(x_i, y_i, z_i)$, l'écart vertical entre la valeur observée $z_i$
et la valeur prédite par le plan $a x_i + b y_i + c$ est
\[
\varepsilon_i = z_i - (a x_i + b y_i + c).
\]

On cherche $(a, b, c)$ qui minimise la somme des carrés de ces écarts :
\[
E(a,b,c) = \sum_{i=1}^{n} \left( z_i - a x_i - b y_i - c \right)^2.
\]

\section{Conditions de minimum}

Le minimum de $E$ est atteint lorsque les dérivées partielles par rapport à
$a$, $b$ et $c$ sont nulles.

\subsection*{Dérivée par rapport à $a$}

\question{Déterminer \frac{\partial E}{\partial a}. Mettre les éléments dépendent de a, b et c à droite du signal égal, les autres à gauche.}

\question{Déterminer \frac{\partial E}{\partial b}. Mettre les éléments dépendent de a, b et c à droite du signal égal, les autres à gauche.}

\question{Déterminer \frac{\partial E}{\partial c}. Mettre les éléments dépendent de a, b et c à droite du signal égal, les autres à gauche.}

\section{Système linéaire (équations normales)}

Les variables suivantes permettent de simplifier l'écriture:
\[
S_x = \sum x_i, \quad
S_y = \sum y_i, \quad
S_{xx} = \sum x_i^2, \quad
S_{yy} = \sum y_i^2, \quad
S_{xy} = \sum x_i y_i,
\]
\[
S_{xz} = \sum x_i z_i, \quad
S_{yz} = \sum y_i z_i, \quad
S_z = \sum z_i,
\]

\question{A partir des trois équations précédentes, déterminer le système linéaire d'inconnues $a, b, c$. Le mettre sous forme matricielle.}

\section{Résolution par pivot de Gauss}

Ce système sera résolu par la méthode du pivot de Gauss, présentée dans le cours d'informatique C15.

\question{Déterminer les valeurs de $a$, $b$ et $c$ grâce à une résolution python par pivot de Gauss.}

\section{Décalage du plan}

Dans le but d'obtenir le plan tangent extérieur matière, la suite va consister à déterminer les points les plus éloignés du plan dans les deux directions afin de décaler le plan.

\question{Déterminer la distance $d_i$ entre un point de coordonnées $(x_i,y_i)$ et le plan ($a$, $b$, $c$).}

\question{Écrire la fonction \verb?distance(point,plan)? qui permet de déterminer $d_i$..}

\question{Trouver l'index $i_{max}$ du point dont la distance positive est la plus éloignée du plan.}

\question{Trouver l'index $i_{min}$ du point dont la distance négative est la plus éloignée du plan.}

\section{Distance d'un point au plan}


\end{document} 

\[
\frac{\partial E}{\partial a}
= \sum_{i=1}^n 2\left( z_i - a x_i - b y_i - c \right)(-x_i) = 0
\]

En développant et en divisant par $-2$ :
\[
\sum_{i=1}^n x_i z_i = a \sum_{i=1}^n x_i^2 + b \sum_{i=1}^n x_i y_i + c \sum_{i=1}^n x_i.
\]

\subsection*{Dérivée par rapport à $b$}

\[
\frac{\partial E}{\partial b}
= \sum_{i=1}^n 2\left( z_i - a x_i - b y_i - c \right)(-y_i) = 0
\]

soit
\[
\sum_{i=1}^n y_i z_i = a \sum_{i=1}^n x_i y_i + b \sum_{i=1}^n y_i^2 + c \sum_{i=1}^n y_i.
\]

\subsection*{Dérivée par rapport à $c$}

\[
\frac{\partial E}{\partial c}
= \sum_{i=1}^n 2\left( z_i - a x_i - b y_i - c \right)(-1) = 0
\]

soit
\[
\sum_{i=1}^n z_i = a \sum_{i=1}^n x_i + b \sum_{i=1}^n y_i + c\, n.
\]

\[
\begin{cases}
S_{xx}\, a + S_{xy}\, b + S_x\, c = S_{xz} \\
S_{xy}\, a + S_{yy}\, b + S_y\, c = S_{yz} \\
S_x\, a    + S_y\, b    + n\, c   = S_z
\end{cases}
\]

ou sous forme matricielle :
\[
\begin{pmatrix}
S_{xx} & S_{xy} & S_x \\
S_{xy} & S_{yy} & S_y \\
S_x    & S_y    & n
\end{pmatrix}
\begin{pmatrix} a \\ b \\ c \end{pmatrix}
=
\begin{pmatrix} S_{xz} \\ S_{yz} \\ S_z \end{pmatrix}.
\]

Ce système $3 \times 3$ est appelé \textbf{système des équations normales}.
Sa matrice est symétrique et, sauf cas dégénéré (points alignés ou tous
identiques en $x$ et $y$), elle est inversible : le système admet une solution
unique $(a, b, c)$.

Le plan $z = ax + by + c$ peut s'écrire sous forme implicite
\[
ax + by - z + c = 0,
\]
de vecteur normal $\vec{n} = (a, b, -1)$.

La distance d'un point $(x_i, y_i, z_i)$ à ce plan est donnée par
\[
d_i = \frac{|a x_i + b y_i - z_i + c|}{\sqrt{a^2 + b^2 + 1}}.
\]

 
